\section{The field \texorpdfstring{$E$}{E} at a prime}\label{sec:field}

This section collects the arithmetic of the CM field $E=K(\zeta)$ of \cref{ss:field-E} that the
divisibility arguments of \cref{part:derivatives} use: a generator of the different, the
decomposition of primes, the Weil data of $E$ and the places chosen by the data. Throughout, $s\ge1$
is fixed, $\zeta\in E$ is the chosen primitive root of unity of order $q_1=q_1(s)$, and $P=\Phi_{q_1}$
denotes the cyclotomic polynomial of order $q_1$; the letter $\Phi$ alone denotes the group of
infinite places. We put
\[
  \beta=\frac{1+\sqrt{-d}}2\in\OO_K,\qquad \beta\bar\beta=\frac{d+1}4,\qquad \beta+\bar\beta=1,
\]
and we write $g_v\in\Gal(E/K)$ for the automorphism attached to the place $v\in\Phi$, so that
$\Phi_v=\Phi_1\circ g_v$. For a finite place $w$ of $E$ and $v\in\Phi$ we write $v\cdot w$ for the
place $g_v^{-1}(w)$, that is, $x\in v\cdot w$ if and only if $g_v(x)\in w$, and $\bar w$ for the
complex conjugate place. The residue field of $w$ has $q_w$ elements.

\subsection{The different}\label{ss:different}

\begin{proposition}\label{prop:different}
Let $z\in\OO_E$ be a primitive root of unity of order $q_1$.
\begin{enumerate}[label=\textup{(\alph*)}]
\item $P'(z)$ divides $q_1$ in $\OO_E$.
\item $\OO_E=\OO_K[z]$, and the different of $E/K$ is the principal ideal $(P'(z))$.
\item $\OO_E=\Z[\beta,\zeta]$, and the different of $E/\Q$ is generated by $P'(\zeta)\sqrt{-d}$. It
contains $q_1d$.
\end{enumerate}
\end{proposition}
\lean{New.Parameters.eval_derivative_cyclotomic_dvd_q₁, New.Parameters.map_differentIdeal_adjoin_eq_span, New.Parameters.span_derivative_cyclotomic_dvd_differentIdeal, New.Parameters.differentIdeal_eq_span_and_adjoin_eq_top, New.Parameters.relDifferent_spec, New.Parameters.adjoin_betaKInE_zetaInt_eq_top, New.Parameters.map_differentIdeal_eq_span_sqrtNegDInE, New.Parameters.sqrtNegDInE_sq, New.Parameters.differentIdeal_eq_span_relDifferentGen_mul, New.Parameters.conductor_mem_differentIdeal}

\begin{proof}
(a) Write $X^{q_1}-1=P(X)Q(X)$ in $\Z[X]$. Differentiating at the root $z$ of $P$ gives
$q_1z^{q_1-1}=P'(z)Q(z)$, and multiplying by $z$ gives $q_1=P'(z)\,zQ(z)$.

(b) The ring of integers of the cyclotomic field $\Q(z)$ is $\Z[z]$, and $P$ is the minimal
polynomial of $z$; by the formula of Euler and Dedekind for monogenic orders, the different of
$\Q(z)/\Q$ is generated by $P'(z)$. The different is transitive in towers, so in $\OO_E$
\[
  \mathfrak D_{E/\Q(z)}\cdot(P'(z))=\mathfrak D_{E/\Q}=\mathfrak D_{E/K}\cdot\mathfrak D_{K/\Q}\OO_E
  =\mathfrak D_{E/K}\cdot(\sqrt{-d}).
\]
The ideal $(P'(z))$ divides $q_1$ by (a), and $q_1$ is prime to $d$, hence to $(\sqrt{-d})$; so
$(P'(z))$ divides $\mathfrak D_{E/K}$. On the other hand $z$ is integral over $\OO_K$ and generates
$E$ over $K$ with minimal polynomial $P$ (\cref{prop:field-E}(a)), so $\mathfrak D_{E/K}$ divides
$P'(z)\,\mathfrak f$, where $\mathfrak f$ is the conductor of $\OO_K[z]$ in $\OO_E$, and in particular
$\mathfrak D_{E/K}\mid(P'(z))$. Hence $\mathfrak D_{E/K}=(P'(z))$, and equality in the Euler--Dedekind
divisibility forces $\mathfrak f=\OO_E$, that is, $\OO_E=\OO_K[z]$.

(c) $\OO_K=\Z[\beta]$ as $d\equiv3\pmod4$, so $\OO_E=\OO_K[\zeta]=\Z[\beta,\zeta]$. The extension
of $\mathfrak D_{K/\Q}=(\sqrt{-d})$ to $\OO_E$ is generated by $\sqrt{-d}$, and transitivity gives
$\mathfrak D_{E/\Q}=(P'(\zeta)\sqrt{-d})$. Finally $P'(\zeta)\mid q_1$ and $\sqrt{-d}\mid d$.
\end{proof}

\begin{proposition}\label{prop:delta-E}
There is $c\in\OO_E$ with $c\bar c=1$ such that $cP'(\zeta)$ is fixed by complex conjugation. The
element
\[
  \delta_E=c\,P'(\zeta)\sqrt{-d}\in\OO_E
\]
generates the different of $E/\Q$, satisfies $\bar\delta_E=-\delta_E$, and
\[
  \prod_{v\in\Phi}\abs{\delta_E}_v=\mathrm{rd}_E^{\,f},\qquad
  \abs{\delta_E}_v=\abs{P'(\zeta)}_v\sqrt d\quad(v\in\Phi).
\]
\end{proposition}
\lean{New.Parameters.relDifferentGen, New.Parameters.relDifferentGen_eq_prod_sub, New.Parameters.exists_unit_mul_zetaInt_sub_real, New.Parameters.exists_unit_mul_relDifferentGen_real, New.Parameters.differentGenPhase, New.Parameters.deltaE, New.Parameters.star_deltaE, New.Parameters.differentIdeal_eq_span_deltaE, New.Parameters.prod_deltaE_eq_rootDiscr_pow, New.Parameters.norm_Φ_differentGenPhase, New.Parameters.infinitePlace_apply_deltaE}

\begin{proof}
Since $P=\prod_a(X-a)$ over the primitive roots of unity $a$ of order $q_1$, we have
$P'(\zeta)=\prod_{a\ne\zeta}(\zeta-a)$. As $q_1$ is even, every such $a$ is $\zeta^{2e+1}$ for some
$e\ge0$. Let $i=\zeta^{q_1/4}$, a square root of $-1$ (recall $4\mid q_1$), and
$u=i\,\bar\zeta^{\,e+1}$. Then $u\bar u=1$ and
\[
  u(\zeta-\zeta^{2e+1})=i\bigl(\zeta^{-e}-\zeta^{e}\bigr)
\]
is the product of the two purely imaginary elements $i$ and $\bar\zeta^{\,e}-\zeta^e$, so it is
real. The product $c$ of these units over $a\neq\zeta$ satisfies $c\bar c=1$, and $cP'(\zeta)$ is
real. Hence $\delta_E$ is real times $\sqrt{-d}$, and $\bar\delta_E=-\delta_E$. It generates the
different by \cref{prop:different}(c), since $c$ is a unit. As $c\bar c=1$, $\abs{c}_v=1$ at every
infinite place, and $\abs{\delta_E}_v=\abs{P'(\zeta)}_v\sqrt d$. Finally $E$ is totally imaginary of
degree $2f$, so $\abs{\Nm_{E/\Q}(\delta_E)}=\prod_v\abs{\delta_E}_v^2$, and the norm of the different
is $\abs{\disc E}=\mathrm{rd}_E^{\,2f}$.
\end{proof}

\begin{remark}\label{rem:delta-size}
At every $v\in\Phi$ one has $\abs{\delta_E}_v=A_v\,T(s)\sqrt d$ with $A_v>0$ and
$\abs{\log A_v}\le C_0$: indeed $P(X)=P_0(X^{T(s)})$ for the cyclotomic polynomial $P_0$ of order
$q_0$, so $P'(\zeta)=T(s)\zeta^{T(s)-1}P_0'(\xi)$ with $\xi=\zeta^{T(s)}$ a primitive root of order
$q_0$, and $\abs{\log\abs{P_0'(\xi)}}$ is one of the terms of $C_0$ in \cref{rem:constant-bound}.
Consequently $\abs{\log\abs{\delta_E}_v-\log\mathrm{rd}_E}\le2C_0$, and for data as in
\cref{def:data} both $\abs{\delta_E}_v$ and $\mathrm{rd}_E$ lie in $[x^{10},x^{12}]$.
\end{remark}
\lean{New.Parameters.exists_abs_log_le_and_apply_deltaE_eq, New.Parameters.abs_log_deltaE_sub_log_rootDiscr_le, New.Data.hasLogBound_deltaE, New.Data.rootDisc_mem_Icc, New.Data.hasLogBound_rootDisc}

\subsection{Decomposition of primes}\label{ss:decomposition}
For a prime $p\nmid q_1d$ the \emph{Frobenius element} at $p$ is the pair
\[
  \mathrm{Fr}_p=\bigl(\chiK(p)p,\ \chiK(p)\bigr)\in\Phi\times\{\pm1\},
\]
where the first entry is the place of the class of $\chiK(p)p$ modulo $q_1$. The group
$\Phi\times\{\pm1\}$ acts on $\OO_E$: $(v,1)$ acts by $g_v$ and $(v,-1)$ by $g_v$ composed with
complex conjugation.
\lean{New.Parameters.frobElem, New.Parameters.IsIdOrConjWithSign, New.Parameters.galOfPlaceInt}

\begin{lemma}\label{lem:frobenius-generators}
Let $p\nmid q_1d$ be prime, $\varphi\colon\OO_E\to R$ a ring homomorphism to a domain of
characteristic $p$, $n\ge0$, and $(c,e)$ either $(\mathrm{id},1)$ or (complex conjugation, $-1$).
For $v\in\Phi$ the following are equivalent:
\begin{enumerate}[label=\textup{(\roman*)}]
\item $\varphi(g_v x)=\varphi(cx)^{p^n}$ for all $x\in\OO_E$;
\item $(v,e)=\mathrm{Fr}_p^{\,n}$.
\end{enumerate}
\end{lemma}
\lean{New.Parameters.ringHom_ext_zetaInt_betaKInE, New.Parameters.forall_map_galOfPlaceInt_eq_pow_iff, New.Parameters.chiK_pow_eq_one_or_neg_one, New.Parameters.map_galOfPlaceInt_zetaInt_inj, New.Parameters.map_galOfPlaceInt_eq_pow_iff_eq_frobElem_pow, New.Parameters.ne_two_of_not_dvd_q₁}

\begin{proof}
Both sides of (i) are ring homomorphisms in $x$, and $\OO_E=\Z[\beta,\zeta]$, so (i) holds if and
only if it holds at $\zeta$ and at $\sqrt{-d}=2\beta-1$ ($p$ is odd, since $2\mid q_1$). At $\zeta$,
$g_v\zeta=\zeta^{a_v}$ with $a_v$ the class of $v$, and $c\zeta=\zeta^e$, so the condition is
$\varphi(\zeta)^{a_v}=\varphi(\zeta)^{ep^n}$. Since $p\nmid q_1$, $\varphi(\zeta)$ is a primitive root
of unity of order $q_1$ in $R$, and the condition says $a_v\equiv ep^n\pmod{q_1}$. At $\sqrt{-d}$, $g_v$
fixes $\sqrt{-d}$, $c\sqrt{-d}=e\sqrt{-d}$, and
$(\sqrt{-d})^{p}=(-d)^{(p-1)/2}\sqrt{-d}\equiv\bigl(\tfrac{-d}p\bigr)\sqrt{-d}=\chiK(p)\sqrt{-d}$ modulo
$p$, by Euler's criterion and quadratic reciprocity ($d\equiv3\pmod4$). Thus the condition is
$\chiK(p)^ne\,\varphi(\sqrt{-d})=\varphi(\sqrt{-d})$. Since $\varphi(\sqrt{-d})^2=-d\ne0$ in $R$ and
$1\neq-1$ in $R$, the condition reads $e=\chiK(p)^n$. The two conditions together say
$(v,e)=\bigl((\chiK(p)p)^n,\chiK(p)^n\bigr)=\mathrm{Fr}_p^{\,n}$.
\end{proof}

\begin{proposition}[Decomposition]\label{prop:decomposition}
Let $p\nmid q_1d$ be prime and $w$ a place of $E$ above $p$.
\begin{enumerate}[label=\textup{(\alph*)}]
\item For $v\in\Phi$: $v\cdot w=w$ if and only if $(v,1)$ is a power of $\mathrm{Fr}_p$, and
$v\cdot w=\bar w$ if and only if $(v,-1)$ is a power of $\mathrm{Fr}_p$.
\item $q_w=p^{o}$, where $o$ is the order of $\mathrm{Fr}_p$; and $p$ is unramified at $w$.
\item Every place of $E$ above $p$ is $v\cdot w$ or $\overline{v\cdot w}$ for some $v\in\Phi$.
\end{enumerate}
\end{proposition}
\lean{New.Parameters.exists_mk_galOfPlaceInt_eq_pow_of_smul_eq, New.Parameters.smul_eq_iff_exists_frobElem_pow, New.Parameters.smul_eq_self_iff_exists_frobElem_pow, New.Parameters.smul_eq_star_iff_exists_frobElem_pow, New.Parameters.q_eq_pow_orderOf_frobElem, New.Parameters.natCast_notMem_sq, New.Parameters.intValuation_natCast_eq_exp_neg_one, New.Parameters.exists_gal_smul_mem_iff, New.Parameters.exists_smul_eq_or_star}

\begin{proof}
(a) Let $c$ be the identity or complex conjugation, with sign $e$, and let $\varphi$ be the reduction
modulo $w$. If $(v,e)=\mathrm{Fr}_p^n$, then by \cref{lem:frobenius-generators} $g_vx\equiv(cx)^{p^n}$
modulo $w$, so $g_vx\in w$ if and only if $cx\in w$; that is, $v\cdot w=c(w)$. Conversely, if
$v\cdot w=c(w)$, then $g_v\circ c^{-1}$ maps $w$ to itself and induces an automorphism of the residue
field $\OO_E/w$, which is a power $y\mapsto y^{p^n}$ of the Frobenius automorphism; so (i) of
\cref{lem:frobenius-generators} holds and $(v,e)=\mathrm{Fr}_p^n$.

(b) The residue degree is the least $n\ge1$ with $y^{p^n}=y$ for all $y\in\OO_E/w$, which by
\cref{lem:frobenius-generators} with $v=1$, $e=1$ is the least $n\ge1$ with $\mathrm{Fr}_p^n=1$. If
$p\in w^2$, then $w$ divides the different of $E/\Q$, which contains $q_1d$
(\cref{prop:different}(c)); then $p\mid q_1d$, which is excluded. Hence $\ord_w(p)=1$.

(c) $\Gal(E/\Q)$ acts transitively on the places above $p$, and every element of $\Gal(E/\Q)$ is some
$g_v$ or some $g_v$ followed by complex conjugation, since $\Gal(E/K)$ has index two and is normal.
\end{proof}

\begin{corollary}[Inert primes]\label{cor:inert}
Let $p$ be prime with $\chiK(p)=-1$ and $w$ a place above $p$. Then $p\nmid q_1d$,
$\mathrm{Fr}_p=(-p,-1)$, and:
\begin{enumerate}[label=\textup{(\alph*)}]
\item $v\cdot w=w$ if and only if $v=(-p)^n$ with $n$ even, and $v\cdot w=\bar w$ if and only if
$v=(-p)^n$ with $n$ odd;
\item $q_w=p^{\operatorname{lcm}(o,2)}$, where $o$ is the order of $-p$ in $(\Z/q_1)^\times$;
\item for odd $n$ and every $y\in\OO_E$, $y^{p^n}\equiv g_{-p^n}(\bar y)\pmod{p\OO_E}$.
\end{enumerate}
\end{corollary}
\lean{New.Parameters.not_dvd_q₁_mul_d_of_inert, New.Parameters.smul_eq_self_iff_of_inert, New.Parameters.smul_eq_star_iff_of_inert, New.Parameters.q_eq_pow_lcm_of_inert, New.Parameters.natCast_dvd_pow_sub_galOfPlaceInt_star}

\begin{proof}
Primes dividing $q_1$ split in $K$ (\cref{lem:split-primes}) and $\chiK(d)=0$, so $p\nmid q_1d$, and
$\mathrm{Fr}_p^n=((-p)^n,(-1)^n)$. Then (a) and (b) are \cref{prop:decomposition}(a),(b): the order of
$\mathrm{Fr}_p$ is $\operatorname{lcm}(o,2)$. For (c) apply \cref{lem:frobenius-generators} to the
reduction modulo each place above $p$, with $v=-p^n$, $c$ complex conjugation and $e=-1=\chiK(p)^n$:
$g_{-p^n}(y)\equiv\bar y^{\,p^n}$ modulo every place above $p$, hence modulo $p$, which is
unramified and so the product of these places. Replacing $y$ by $\bar y$ gives (c).
\end{proof}

\begin{proposition}[Primes $\equiv1$ modulo $q_1d$]\label{prop:split-one}
Let $p\equiv1\pmod{q_1d}$ be prime and $w$ a place above $p$. Then $\chiK(p)=1$,
$\mathrm{Fr}_p=1$, $q_w=p$, $v\cdot w=w$ only for $v=1$, and $v\cdot w\ne\bar w$ for all $v$. If
$a\in\OO_E$ satisfies $a\bar a=pc$ with $c\in\N$ prime to $p$, then $w$ contains exactly one of $a$
and $\bar a$, with $\ord_w(a)=1$ if $a\in w$, and $v\mapsto$ (the one of $v\cdot w$,
$\overline{v\cdot w}$ that contains $a$) is a bijection from $\Phi$ onto the places above $p$ that
contain $a$.
\end{proposition}
\lean{New.Parameters.chiK_eq_one_of_modEq_one, New.Parameters.not_dvd_q₁_mul_d_of_modEq_one, New.Parameters.frobElem_eq_one, New.Parameters.smul_ne_star, New.Parameters.smul_eq_self_iff_eq_one, New.Parameters.q_eq_of_modEq_one, New.Parameters.star_mem_iff_notMem, New.Parameters.intValuation_eq_of_mul_star, New.Parameters.eq_of_smul_eq_or_smul_eq_star, New.Parameters.nonempty_infinitePlace_equiv_finPlace}

\begin{proof}
$\chiK(p)=\chiK(1)=1$ since $\chiK$ has conductor $d$, and $p\equiv1\pmod{q_1}$, so
$\mathrm{Fr}_p=(1,1)$. \Cref{prop:decomposition} gives the statements on $v\cdot w$ and $q_w=p$.
If $a$ and $\bar a$ both lay in $w$, then $pc\in w^2$; but $\ord_w(p)=1$ and $c$ is a unit at $w$.
If neither did, $pc\notin w$. Hence exactly one lies in $w$, and
$\ord_w(a)+\ord_w(\bar a)=\ord_w(pc)=1$. For the bijection: the map is defined, since exactly one of
$v\cdot w$ and its conjugate contains $a$ (apply the previous statement to the place $v\cdot w$).
It is injective, since $v\cdot w\in\{u\cdot w,\overline{u\cdot w}\}$ forces $u^{-1}v$ to fix $w$ or to
conjugate it, hence $u=v$. It is surjective by \cref{prop:decomposition}(c).
\end{proof}

An automorphism $\sigma$ of $\C$ over $K$ permutes the embeddings $\Phi_v$: $\sigma\circ\Phi_v=
\Phi_{\sigma\cdot v}$, where $\sigma\cdot v$ is the product in $\Phi$ of $v$ and the place of the
embedding $\sigma\circ\Phi_1$. We write $v\mapsto\sigma\cdot v$ for this permutation of $\Phi$.
\lean{New.Parameters.placePerm, New.Parameters.placePerm_apply, New.Parameters.apply_Phi_eq_Phi_placePerm}

\subsection{The Weil data of \texorpdfstring{$E$}{E}}\label{ss:weil-data-E}
The Weil representation of a hermitian space over $E$ (\cref{ss:weil-rep}) depends on \emph{Weil
data}: a purely imaginary $\jmath\in E^\times$, a finite idele $\alpha$ of $E$, and a unitary Hecke
character $\upsilon$ of $E$ (a \emph{splitting character}) whose restriction to $K$ is the quadratic
character attached to $E/K$ times the prescribed archimedean factors. We now fix them.

\begin{definition}\label{def:alpha}
The finite idele $\alpha$ of $E$ has component $P'(\zeta)^{-1}$ at the places $w$ with $\beta\in w$
and $1$ at all other places. The lattice $\mathfrak a_0=E\cap\alpha\widehat\OO_E$ is the set of
$y\in E$ with $\alpha^{-1}y$ integral at all finite places.
\end{definition}
\lean{New.Parameters.α, New.Parameters.alpha_apply, New.Parameters.inv_alpha_apply, New.diagonalAtPlaces, New.diagonalAtPlaces_apply}

\begin{lemma}\label{lem:alpha}
\begin{enumerate}[label=\textup{(\alph*)}]
\item A place $w$ with $P'(\zeta)\in w$ contains $\beta$ or $\bar\beta$; and $\ord_w P'(\zeta)=
\ord_{\bar w}P'(\zeta)$.
\item $\alpha^{-1}$ and $q_1\alpha$ are integral; $\alpha\bar\alpha P'(\zeta)$ is a unit at every
finite place.
\item With $\jmath=-1/\sqrt{-d}$, the idele $\jmath\alpha\bar\alpha$ has order
$-\ord_w\mathfrak D_{E/\Q}$ at every $w$; that is, $\jmath\,\mathfrak a_0\bar{\mathfrak a}_0$ is the
inverse different of $E$. Moreover $\Phi_\tau(\jmath)=i/\sqrt d$ for every $\tau$, and
$(1-\jmath^{-1})/2=\beta$.
\item Every root of unity $y\in E$ is $\delta/\bar\delta$ for a unit $\delta\in\OO_E^\times$ with
$\prod_{v}\Phi_v(\delta)=1$.
\end{enumerate}
\end{lemma}
\lean{New.Parameters.valuation_map_relDifferentGen, New.Parameters.valuation_star_relDifferentGen, New.Parameters.betaKInE_mem_or_star_mem, New.Parameters.inv_alpha_mem, New.Parameters.v_alpha_apply, New.Parameters.v_star_alpha_apply, New.Parameters.v_alpha_mul_star_mul_relDifferentGen, New.Parameters.q₁_mul_alpha_mem, New.Parameters.v_neg_inv_sqrtNegDInE_mul_alpha_mul_star, New.Parameters.v_weilData_j_mul_alpha_mul_star, New.Parameters.weilData_j_eq, New.Parameters.phi_weilData_j, New.Parameters.traceOne_eq_betaKInE, New.Parameters.norm_zeta_pow_eq_one, New.Parameters.norm_zeta_sub_one_eq_one, New.Parameters.exists_unit_prod_phi_eq_one_and_mul_star_eq}

\begin{proof}
(a) $P'(\zeta)$ divides $q_1$, all of whose prime factors divide $\beta\bar\beta=(d+1)/4$; so a prime
$w$ containing $P'(\zeta)$ contains $\beta$ or $\bar\beta$. Any injective endomorphism of $\OO_E$
sends $\zeta$ to a primitive root of order $q_1$, so it sends $P'(\zeta)$ to another generator of
$\mathfrak D_{E/K}$ (\cref{prop:different}(b)); applied to complex conjugation this gives
$\ord_{\bar w}P'(\zeta)=\ord_wP'(\zeta)$.

(b) $\alpha^{-1}$ has components $P'(\zeta)$ or $1$, which are integral; and $q_1\alpha$ has
components $q_1/P'(\zeta)$ or $q_1$, integral by \cref{prop:different}(a). No place contains both
$\beta$ and $\bar\beta$, since $\beta+\bar\beta=1$. If $\beta\in w$, then $\alpha_w=P'(\zeta)^{-1}$
and $\bar\alpha_w=1$ (because $\bar\alpha_w$ is the conjugate of $\alpha_{\bar w}$ and $\beta\notin\bar
w$); if $\bar\beta\in w$, then $\alpha_w=1$ and $\bar\alpha_w$ has the order of $P'(\zeta)^{-1}$ at
$w$ by (a); otherwise $\alpha_w=\bar\alpha_w=1$ and $P'(\zeta)$ is a unit at $w$ by (a).

(c) By (b), $\jmath\alpha\bar\alpha$ has the order of $(P'(\zeta)\sqrt{-d})^{-1}$, which is the order of
$\mathfrak D_{E/\Q}^{-1}$ by \cref{prop:different}(c). The embedding $\Phi_\tau$ is over $K\subset\C$,
so $\Phi_\tau(\sqrt{-d})=i\sqrt d$ and $\Phi_\tau(\jmath)=-1/(i\sqrt d)=i/\sqrt d$; and
$(1-\jmath^{-1})/2=(1+\sqrt{-d})/2$.

(d) Let $q_1=4m$ and $i=\zeta^m$. The conjugates of $\zeta$ over $K$ are the $\zeta^a$ with $a$ a
unit modulo $q_1$, so $\Nm_{E/K}(\zeta^k)=\zeta^{k\sum_aa}=1$, because $a\mapsto-a$ pairs the units and
$\sum_aa\equiv0\pmod{q_1}$. Likewise $\Nm_{E/K}(\zeta-1)=\prod_a(\zeta^a-1)=P(1)$, as $\totient(q_1)$ is
even, and $P(1)=1$ because $q_1$, divisible by $2$ and by $3$, is not a prime power; so $\zeta-1$ is
a unit of norm one. Put $\eta=i(\zeta-1)$; then
\[
  \frac{\eta}{\bar\eta}=\frac{i(\zeta-1)}{-i(\zeta^{-1}-1)}=-\frac{\zeta-1}{-\zeta^{-1}(\zeta-1)}=\zeta .
\]
By \cref{prop:field-E}(d) every root of unity is $\zeta^k$, and $\delta=\eta^k$ serves; by
\cref{prop:field-E}(e) the product of the $\Phi_v(\delta)$ is $\Nm_{E/K}(\delta)=1$.
\end{proof}

\begin{proposition}\label{prop:weil-data-E}
The element $\jmath=-1/\sqrt{-d}$, the idele $\alpha$ of \cref{def:alpha} and a splitting character
$\upsilon$ are Weil data for $E$; such a character exists. We fix it.
\end{proposition}
\lean{New.Parameters.instWeilDataE}

\begin{proof}
The conditions on $\jmath$ and $\alpha$ are that $\jmath$ is purely imaginary and that
$\jmath\alpha\bar\alpha$ generates the inverse different locally at every finite place; both hold by
\cref{lem:alpha}(c). A splitting character exists as soon as every root of unity of $E$ has the form
$\delta/\bar\delta$ with $\delta$ a unit of norm one to $K$; this is the criterion of
\cref{thm:B-splitting-character}, and the hypothesis is \cref{lem:alpha}(d).
\end{proof}

\begin{definition}\label{def:vacuum-root}
Let $\chi_{\mathrm f}(x)=\upsilon(x)\abs x^{-1/2}$ for finite ideles $x$ (extended by $1$ at the
infinite places). The \emph{vacuum root} $s_0\in\C$ is the principal $f$-th root of
$\chi_{\mathrm f}(\alpha)^{-1}=\upsilon(\alpha)^{-1}\abs{\alpha}^{1/2}$.
\end{definition}
\lean{New.Parameters.vacuumRoot, New.Parameters.vacuumRoot_pow_numInfPlaces}

\begin{proposition}\label{prop:vacuum-root}
A positive power of $s_0$ lies in $K$, and for every automorphism $\sigma$ of $\C$ over $K$
\[
  \abs{\sigma(s_0)}^2=\frac{\mathrm{rd}_E}{\sqrt d}.
\]
\end{proposition}
\lean{New.Parameters.exists_chiFin_alpha_pow_mem, New.Parameters.exists_vacuumRoot_pow_mem, New.Parameters.ideleAbs_alpha_eq_prod, New.Parameters.norm_gal_vacuumRoot}

\begin{proof}
A positive power of $\chi_{\mathrm f}(\alpha)$ lies in $K$: $\upsilon$ has finite order on the unit
ideles, and a power of $\alpha$ is, up to such units, the idele of a principal ideal $(y)$, $y\in E$,
at which the value of $\upsilon$ is the archimedean factor, a norm from $E$ to $K$
(\cref{thm:B-splitting-character}). Hence $s_0^n\in K$ for some $n\ge1$, and $\abs{\sigma s_0}^n=
\abs{\sigma(s_0^n)}=\abs{s_0^n}$, so $\abs{\sigma s_0}=\abs{s_0}$. Next, $\abs{s_0}^{2f}=\abs\alpha$ as
$\upsilon$ is unitary. Since $\alpha\bar\alpha P'(\zeta)$ is a unit at every finite place
(\cref{lem:alpha}(b)) and $\abs{\bar\alpha}=\abs\alpha$, we have
$\abs\alpha^2\,\abs{P'(\zeta)}_{\mathrm f}=1$, and the product formula gives
$\abs{P'(\zeta)}_{\mathrm f}^{-1}=\prod_v\abs{P'(\zeta)}_v^2$. Hence
$\abs{s_0}^{2f}=\prod_v\abs{P'(\zeta)}_v=\prod_v\abs{\delta_E}_v/\sqrt d^{\,f}=(\mathrm{rd}_E/\sqrt d)^f$ by
\cref{prop:delta-E}.
\end{proof}

\subsection{The chosen places and characters of given weight}\label{ss:chosen}
For the rest of \cref{sec:field} the data of \cref{def:data} are fixed. Recall that $M=p_0p_1=\pi_M\bar\pi_M$
with $\bar\pi_M\in\OO_K$. The \emph{chosen places} are the finite places $w$ of $E$ with $\pi_M\in w$.
\lean{New.Data.piM, New.Data.IsChosen, New.Data.ChosenPlace, New.Data.piM_mul_star}

\begin{lemma}\label{lem:chosen}
Every chosen place lies above $p_0$ or $p_1$ and has norm $q_w=p_i$; there are finitely many. For
$i=0,1$, $\pi_M\bar\pi_M=p_i c_i$ with $c_i\in\N$ prime to $p_i$, so by \cref{prop:split-one} the
chosen places above $p_i$ are in bijection with $\Phi$.
\end{lemma}
\lean{New.Data.piMBar_ne_zero, New.Data.piM_ne_zero, New.Data.exists_piM_mul_star_eq_levelPrime_mul_coprime,New.Data.exists_levelPrime_mem_of_isChosen, New.Data.isChosen_iff, New.Data.q_eq_levelPrime, New.Data.finite_isChosen}

\begin{proof}
A chosen $w$ contains $\pi_M\bar\pi_M=p_0p_1$, hence $p_0$ or $p_1$. The primes $p_i$ are $\equiv1$
modulo $q_2^2$, hence modulo $q_1d$ (\cref{lem:split-primes}), so \cref{prop:split-one} applies with
$a=\pi_M$ and $c_i=p_{1-i}$.
\end{proof}

At a chosen place $w$ the residue field has $p_i$ elements and $q_2^2\mid p_i-1$, so it carries a
multiplicative character of order $q_2^2$; we fix one, $\omega_w^\flat$, and put
$\omega_w=(\omega^\flat_w)^{q_2}$, of order $q_2$. At the other places $\omega_w=1$. Since $q_1\mid
q_2$, $\omega_w(y)=1$ for every $y$ in the residue field with $y^{q_1}=1$.
\lean{New.Data.exists_mulChar_orderOf_eq_sq_of_isChosen, New.Data.chiMLocalRoot, New.Data.orderOf_chiMLocalRoot, New.Data.chiMLocal, New.Data.chiMLocal_of_isChosen, New.Data.orderOf_chiMLocal, New.Data.chiMLocal_eq_one_of_not_isChosen, New.Data.chiMLocal_apply_eq_one_of_pow_q₁_eq_one, New.Data.chiMLocal_zpow_mul_pow_eq_one}

Let $\U(1)(\Ad)$ be the group of ideles $u$ of $E$ with $u\bar u=1$. A continuous character $\mu$ of
$\U(1)(\Ad)$, trivial on $\U(1)(E)$, \emph{has weight $e\colon\Phi\to\Z$ and tame part
$\psi=(\psi_w)_w$} if almost all $\psi_w$ are trivial and
\[
  \mu(u)=\prod_{v\in\Phi}u_v^{\,e(v)}\prod_{w}\psi_w(u\bmod w)
\]
for every $u$ with integral finite part, where $u_v=\Phi_v(u)$. Let
$w_0(v)=\kappa_0(v)-\Delta/2=\kappa(v)+r^8$ be the lift weight $\kappa_0$ of \cref{def:data} shifted
by half the gap.
\lean{New.AutChar, New.AutChar.HasWeight}

\begin{proposition}\label{prop:characters-of-weight}
For all $a,n\in\N$ there is a character of $\U(1)(\Ad)$, trivial on $\U(1)(E)$, of weight
$-a\,w_0$ and tame part $(\omega_w^n)_w$. We fix such characters $\chi_0$ (for $a=1$, $n=0$) and
$\chi_M$ (for $a=0$, $n=1$).
\end{proposition}
\lean{New.Data.exists_zeta_pow_eq_of_isIntegral, New.Data.prod_phi_zeta_pow_pow_liftWeight_zero_eq_one, New.Data.exists_char_hasWeight, New.Data.chiZero, New.Data.chiM, New.Data.chiZero_hasWeight, New.Data.chiM_hasWeight, New.AutChar.exists_hasWeight_of_prod_eq_one}

\begin{proof}
By \cref{thm:B-characters-of-weight} such a character exists if the formula for $\mu(u)$ gives $1$
at every $\gamma\in\U(1)(E)$ that is integral. Such $\gamma$ is an algebraic integer all of whose
conjugates have absolute value $1$, hence a root of unity, hence $\gamma=\zeta^k$
(\cref{prop:field-E}(d)). At a finite place the residue of $\zeta^k$ is a $q_1$-th root of unity, so
$\omega_w(\zeta^k)=1$. At the infinite places,
$\prod_v\Phi_v(\zeta^k)^{w_0(v)}=1$: the weight $\kappa(v)$ is a multiple of $q_1$ (both $b$ and
$a/x_1$ are, since $q_1x_1\mid b$ and $a=xb$), so $w_0(v)\equiv r^8\pmod{q_1}$ for all $v$, and
$\prod_v\Phi_v(\zeta^k)^{r^8}=\Nm_{E/K}(\zeta^k)^{r^8}=1$.
\end{proof}

\begin{definition}\label{def:critical-pair}
A \emph{pair} for the data is an integer $z$ with $1\le z\le t$ together with a character $\nu$ of
$\U(1)(\Ad)$, trivial on $\U(1)(E)$, of weight $-c$ with $z\le c(v)\le4b$ for all $v\in\Phi$ and
tame part trivial off the chosen places.
\end{definition}
\lean{New.Data.CriticalPair, New.Data.CriticalPair.le_negWeight}

\begin{proposition}\label{prop:finite-pairs}
There are finitely many pairs.
\end{proposition}
\lean{New.Data.finite_criticalPair, New.AutChar.finite_setOf_hasWeight}

\begin{proof}
The weight $c$ and the integer $z$ range over finite sets; the tame part is a family of characters of
the residue fields at the finitely many chosen places; and for a given weight and tame part there are
finitely many characters, since two of them differ by a character of the compact group
$\U(1)(E)\backslash\U(1)(\Ad)$ that is trivial on the open subgroup of $u$ with $u_v=1$ for $v\in\Phi$
and $u_w\equiv1$ modulo $w$ for all $w$, a subgroup of finite index
(\cref{thm:B-characters-of-weight}).
\end{proof}
